%!TEX root = ../../main.tex
%% ===============================================================
%% 1.2 문제 정의와 도전 과제
%% 파일: chapters/01_introduction/02_challenges.tex
%% 상위: chapters/01_introduction/chapter.tex
%% 정의 라벨: sec:intro-challenges
%% 참조 라벨: -
%% 포함 객체: -
%% 인용: berkson1946limitations, quinonero2009dataset, shimodaira2000improving
%% 상태: 초안 (docs/OUTLINE.md 참고)
%% ===============================================================

\section{문제 정의와 도전 과제}
\label{sec:intro-challenges}

본 논문이 다루는 문제는 다음과 같이 요약된다. \emph{공개 텔레메트리만으로} (i) 경기 안에서 교전이 언제 어디서 일어났는지를 사람의 주석 없이 재현 가능하게 찾아내고(교전 국소화, localization), (ii) 교전이 시작되기 직전 30 초의 관측만으로 어느 팀이 교전에서 이길지를 예측한다(교전 결과 예측). 이 문제는 다음 다섯 가지 이유로 어렵다.

\begin{enumerate}[leftmargin=2em]
  \item \textbf{시공간 복잡성.} 열 명의 플레이어가 동시에 이동하며, 누가 누구와 가까운지(상호작용 그래프)가 연속적으로 바뀐다. 교전의 발생과 결과 모두 이 동적 그래프의 함수이다.
  \item \textbf{이질적 모달리티.} 예측 신호는 플레이어별 수치 상태(골드, 레벨, 스탯), 팀 단위 집계(오브젝트 차이), 이산 사건(킬, 와드), 범주형 메타데이터(챔피언, 룬, 소환사 주문)에 흩어져 있다.
  \item \textbf{시간 해상도 불일치.} 상태는 60 초, 사건은 밀리초 단위이다. 두 시간 축을 정렬할 때 \emph{미래 정보}가 관측 창에 새어 들어가지 않도록 하는 계약(contract)이 필요하다.
  \item \textbf{패치에 따른 공변량 이동.} LoL은 약 2 주마다 밸런스 패치를 받는다. 학습 분포와 시험 분포가 달라지는 공변량 이동(covariate shift) \cite{shimodaira2000improving,quinonero2009dataset}이 구조적으로 존재한다.
  \item \textbf{선택 편향.} 교전을 ``예측 지평 안에 사건이 있는 창''으로 정의하면, 결과에 조건화된 표본 선택이 일어나 Berkson의 역설 \cite{berkson1946limitations}과 같은 편향이 학습 분포에 유입된다.
\end{enumerate}
